% GENERATED FILE. Edit canonical lessons, then run npm run generate:books.
% canonical-source-hash: fnv1a64:b599795bc5608590
% canonical-lessons: JA-C218-nyuuin, JA-C218-taiin, JA-C218-shujutsu, JA-C218-kizu, JA-C218-inochi

\chapter{A Stay in Hospital}
\label{ch:ja-a-stay-in-hospital}
\hlchaptermodality{218}

\begin{chapteropening}
\textbf{By the end of this chapter:} \emph{I can ask about a hospital stay and an operation.}

Complete the last lesson of chapter 218: I can ask about a hospital stay and an operation.
\end{chapteropening}

\section[\texorpdfstring{\ja{にゅういん} (nyūin)}{にゅういん (nyūin)}]{\ja{にゅういん} (nyūin) — going into hospital; a hospital stay}

\label{lesson:JA-C218-nyuuin}



\begin{quote}
Before the new one: say the Japanese for a dentist, then the Japanese for a cavity.
\end{quote}

\subsection*{You'll want to know: \ja{にゅういん}}
\textbf{\ja{にゅういん}} — \emph{nyūin} — \textquotedblleft{}going into hospital; a hospital stay\textquotedblright{}.

\begin{quote}
\textbf{\ja{にゅういんは} \ja{ながいですか。}} — \emph{nyūin wa nagai desu ka} — \textquotedblleft{}Is the hospital stay long?\textquotedblright{}
\end{quote}

\begin{cousinweb}
\textbf{\ja{にゅういん}} (\emph{nyūin}) is Sino-Japanese (kango): two Chinese characters read with their on'yomi, readings borrowed from Chinese. They are \emph{nyū}, \textquotedblleft{}to enter\textquotedblright{} (Middle Chinese \emph{nyip}) and \emph{in}, \textquotedblleft{}courtyard, institution\textquotedblright{} (\emph{hjwenH}) (KANJIDIC2; Middle Chinese in Baxter's transcription).
\end{cousinweb}

\begin{grammarlens}[title={What you've built}]
One more word for everyday life at A2.
\end{grammarlens}

\subsection*{Your turn}
\begin{itemize}
\raggedright
  \item \emph{Say it:} \emph{nyūin}
  \item \emph{Say it:} \emph{nyūin}, once more
  \item \emph{Say it:} \emph{nyūin wa nagai desu ka} — Is the hospital stay long?
  \item \emph{Recall:} say the Japanese for a dentist, then the Japanese for a cavity, then say \emph{nyūin} again
\end{itemize}

\begin{tcolorbox}[breakable,colback=teal!4,colframe=teal!35!black,title={Before you move on}]
What does \ja{にゅういん} mean? (\textquotedblleft{}Going into hospital; a hospital stay\textquotedblright{}.) Say it once more.
\end{tcolorbox}
\section[\texorpdfstring{\ja{たいいん} (taiin)}{たいいん (taiin)}]{\ja{たいいん} (taiin) — leaving hospital}

\label{lesson:JA-C218-taiin}



\begin{quote}
Before the new one: say the Japanese for a cavity, then the Japanese for going into hospital.
\end{quote}

\subsection*{You'll want to know: \ja{たいいん}}
\textbf{\ja{たいいん}} — \emph{taiin} — \textquotedblleft{}leaving hospital\textquotedblright{}.

\begin{quote}
\textbf{\ja{たいいんは} \ja{あしたですか。}} — \emph{taiin wa ashita desu ka} — \textquotedblleft{}Do you leave hospital tomorrow?\textquotedblright{}
\end{quote}

\begin{cousinweb}
\textbf{\ja{たいいん}} (\emph{taiin}) is Sino-Japanese (kango): two Chinese characters read with their on'yomi, readings borrowed from Chinese. They are \emph{tai}, \textquotedblleft{}to withdraw\textquotedblright{} (Middle Chinese \emph{thwojH}) and \emph{in}, \textquotedblleft{}courtyard, institution\textquotedblright{} (\emph{hjwenH}) (KANJIDIC2; Middle Chinese in Baxter's transcription).
\end{cousinweb}

\begin{grammarlens}[title={What you've built}]
One more word for everyday life at A2.
\end{grammarlens}

\subsection*{Your turn}
\begin{itemize}
\raggedright
  \item \emph{Say it:} \emph{taiin}
  \item \emph{Say it:} \emph{taiin}, once more
  \item \emph{Say it:} \emph{taiin wa ashita desu ka} — Do you leave hospital tomorrow?
  \item \emph{Recall:} say the Japanese for a cavity, then the Japanese for going into hospital, then say \emph{taiin} again
\end{itemize}

\begin{tcolorbox}[breakable,colback=teal!4,colframe=teal!35!black,title={Before you move on}]
What does \ja{たいいん} mean? (\textquotedblleft{}Leaving hospital\textquotedblright{}.) Say it once more.
\end{tcolorbox}
\section[\texorpdfstring{\ja{しゅじゅつ} (shujutsu)}{しゅじゅつ (shujutsu)}]{\ja{しゅじゅつ} (shujutsu) — surgery; an operation}

\label{lesson:JA-C218-shujutsu}



\begin{quote}
Before the new one: say the Japanese for going into hospital, then the Japanese for leaving hospital.
\end{quote}

\subsection*{You'll want to know: \ja{しゅじゅつ}}
\textbf{\ja{しゅじゅつ}} — \emph{shujutsu} — \textquotedblleft{}surgery; an operation\textquotedblright{}.

\begin{quote}
\textbf{\ja{しゅじゅつは} \ja{ひつようですか。}} — \emph{shujutsu wa hitsuyō desu ka} — \textquotedblleft{}Is surgery needed?\textquotedblright{}
\end{quote}

\begin{cousinweb}
\textbf{\ja{しゅじゅつ}} (\emph{shujutsu}) is Sino-Japanese (kango): two Chinese characters read with their on'yomi, readings borrowed from Chinese. They are \emph{shu}, \textquotedblleft{}hand\textquotedblright{} (Middle Chinese \emph{syuwX}) and \emph{jutsu}, \textquotedblleft{}skill\textquotedblright{} (\emph{zywit}) (KANJIDIC2; Middle Chinese in Baxter's transcription).
\end{cousinweb}

\begin{grammarlens}[title={What you've built}]
One more word for everyday life at A2.
\end{grammarlens}

\subsection*{Your turn}
\begin{itemize}
\raggedright
  \item \emph{Say it:} \emph{shujutsu}
  \item \emph{Say it:} \emph{shujutsu}, once more
  \item \emph{Say it:} \emph{shujutsu wa hitsuyō desu ka} — Is surgery needed?
  \item \emph{Recall:} say the Japanese for going into hospital, then the Japanese for leaving hospital, then say \emph{shujutsu} again
\end{itemize}

\begin{tcolorbox}[breakable,colback=teal!4,colframe=teal!35!black,title={Before you move on}]
What does \ja{しゅじゅつ} mean? (\textquotedblleft{}Surgery; an operation\textquotedblright{}.) Say it once more.
\end{tcolorbox}
\section[\texorpdfstring{\ja{きず} (kizu)}{きず (kizu)}]{\ja{きず} (kizu) — a wound; a cut}

\label{lesson:JA-C218-kizu}



\begin{quote}
Before the new one: say the Japanese for leaving hospital, then the Japanese for surgery.
\end{quote}

\subsection*{You'll want to know: \ja{きず}}
\textbf{\ja{きず}} — \emph{kizu} — \textquotedblleft{}a wound; a cut\textquotedblright{}.

\begin{quote}
\textbf{\ja{きずは} \ja{ふかいですか。}} — \emph{kizu wa fukai desu ka} — \textquotedblleft{}Is the wound deep?\textquotedblright{}
\end{quote}

\begin{cousinweb}
\textbf{\ja{きず}} (\emph{kizu}) is a native Japanese word (wago), inherited rather than borrowed from Chinese; its kanji was attached later as a kun'yomi, and its older origin is not settled (JMdict, KANJIDIC2).
\end{cousinweb}

\begin{grammarlens}[title={What you've built}]
One more word for everyday life at A2.
\end{grammarlens}

\subsection*{Your turn}
\begin{itemize}
\raggedright
  \item \emph{Say it:} \emph{kizu}
  \item \emph{Say it:} \emph{kizu}, once more
  \item \emph{Say it:} \emph{kizu wa fukai desu ka} — Is the wound deep?
  \item \emph{Recall:} say the Japanese for leaving hospital, then the Japanese for surgery, then say \emph{kizu} again
\end{itemize}

\begin{tcolorbox}[breakable,colback=teal!4,colframe=teal!35!black,title={Before you move on}]
What does \ja{きず} mean? (\textquotedblleft{}A wound; a cut\textquotedblright{}.) Say it once more.
\end{tcolorbox}
\section[\texorpdfstring{\ja{いのち} (inochi)}{いのち (inochi)}]{\ja{いのち} (inochi) — life}

\label{lesson:JA-C218-inochi}



\begin{quote}
Before the new one: say the Japanese for surgery, then the Japanese for a wound.
\end{quote}

\subsection*{You'll want to know: \ja{いのち}}
\textbf{\ja{いのち}} — \emph{inochi} — \textquotedblleft{}life\textquotedblright{}.

\begin{quote}
\textbf{\ja{いのちは} \ja{たいせつです。}} — \emph{inochi wa taisetsu desu} — \textquotedblleft{}Life is precious.\textquotedblright{}
\end{quote}

\begin{cousinweb}
\textbf{\ja{いのち}} (\emph{inochi}) is a native Japanese word (wago), inherited rather than borrowed from Chinese; its kanji was attached later as a kun'yomi, and its older origin is not settled (JMdict, KANJIDIC2).
\end{cousinweb}

\begin{grammarlens}[title={What you've built}]
One more word for everyday life at A2.
\end{grammarlens}

\subsection*{Your turn}
\begin{itemize}
\raggedright
  \item \emph{Say it:} \emph{inochi}
  \item \emph{Say it:} \emph{inochi}, once more
  \item \emph{Say it:} \emph{inochi wa taisetsu desu} — Life is precious.
  \item \emph{Recall:} say the Japanese for surgery, then the Japanese for a wound, then say \emph{inochi} again
\end{itemize}

\begin{tcolorbox}[breakable,colback=teal!4,colframe=teal!35!black,title={Before you move on}]
What does \ja{いのち} mean? (\textquotedblleft{}Life\textquotedblright{}.) Say it once more.
\end{tcolorbox}
